\setcounter{section}{1}
\setcounter{theorem}{7}
	\section{Preliminaries}\label{sec:preliminaries}
	
	\subsection{Notation}
	
	Let $n\in\mathbb{N}$; we denote $\brkts{1,\dots,n}$ by $[n]$, and denote by $\uniform_n$ the uniform distribution over $\bool^n$. We use $\widetilde{O}(\cdot)$ (and $\widetilde{\Omega}(\cdot)$)  to hide polylogarithmic factors. That is, for any $f\colon\mathbb{N}\to\mathbb{N}$, we have that $\widetilde{O}(f(n))=O\cprn{f(n)\cdot\polylog\cprn{f(n)}}$.
	
	In this paper, we will mainly use $\boolpm$ as the Boolean basis. In some parts of this paper, we will use the $\bool$ basis for the simplicity of the presentation. This will be specified in corresponding sections.
	
	\subsection{De Morgan formulas and extensions}
	
	\begin{definition}
	    An $n$-variate \emph{de Morgan formula} is a directed rooted tree; its non-leaf vertices (henceforth, \emph{internal gates}) take labels from $\brkts{{\rm AND},{\rm OR},{\rm NOT}}=\brkts{\land,\lor,\neg}$ and its leaves (henceforth, \emph{variable gates}) take labels from the set of variables $\brkts{x_1,\dots,x_n}$. Each internal gate has bounded in-degree (henceforth, \emph{fan-in}); the {\rm NOT} gate in particular has fan-in $1$ and every variable gate has fan-in $0$. The \emph{size} of a de Morgan formula is the number of its leaf gates.
%
	\end{definition}
	
	In this work, we denote by $\FORMULA[s]$ the class of Boolean functions computable by size-$s$ de Morgan formulas. Let $\mathcal{G}$ denote some class of Boolean functions; then, we denote by $\FORMULA[s]\circ\mathcal{G}$ the class of functions computable by some size-$s$ de Morgan formula where its leaves are labelled by functions in $\mathcal{G}$.
	
%
	
%
%
%
%
%
%
%
%
%
    
%
%
%
%
%
    
    \subsection{Approximating polynomials}
    
    \begin{definition}[Point-wise approximation]
        For a Boolean function $f\colon\boolpm^n\to\boolpm$, we say that the function $\tilde{f}\colon\boolpm^n\to\real$ $\varepsilon$-approximates $f$ if for every $z\in\boolpm^n$,
		\[
			\left| f(z)-\tilde{f}(z)\right| \leq \varepsilon.
		\]
    \end{definition}
    
    We will need the following powerful result for the approximating degree of de Morgan formulas.
    
    \begin{theorem}[\cite{Rei11b}, see also~\cite{ BNRdW07}]\label{thm:approx-polynomial}
		Let $s>0$ be an integer and $0<\varepsilon<1$. Any de Morgan formula $F\colon\boolpm^n\to\boolpm$ of size $s$ has a $\varepsilon$-approximating polynomial of degree $d=O(\sqrt{s}\cdot\log(1/\varepsilon))$. That is, there exists a degree-$d$ polynomial $p\colon\boolpm^n\to\real$ over the reals such that for every $z\in\boolpm^n$,
		\[
			\left| p(z)-F(z)\right| \leq \varepsilon.
		\]
	\end{theorem}
	
	Note that \Cref{thm:approx-polynomial} still holds if we use $\bool$ as the Boolean basis.
    
	\subsection{Communication complexity}
	
	We use standard definitions from communication complexity. In this paper we consider the standard two party model of Yao and its generalizations to multiparty setting. We denote deterministic communication complexity of a Boolean function by $D(f)$ in the two party setting. We refer to \citep{Kushilevitz-Nisan97} for 
	standard definitions from communication complexity.
	
%
%
	    
%
%
	
%
	
	\begin{definition}
	    Let $f\colon\bool^{n}\to\bool$ be a Boolean function. The \emph{communication matrix of $f$}, namely $M_f$, is a $2^{n/2}\times2^{n/2}$ matrix defined by $\prn{M_f}_{x,y}:=f(x,y)$.
	\end{definition}
	
	\begin{definition}
	    A \emph{rectangle} is a set of the form $A\times B$, for $A,B\subseteq\bool^n$. A \emph{monochromatic rectangle} is a rectangle $S$ such that for all pairs $(x,y)\in S$ the value $f(x,y)$ is the same.
    \end{definition}
	
	\begin{lemma}
	    Let $\Pi$ be a protocol that computes $f\colon\bool^{n}\to\bool$ with at most $D$ bits of communication. Then, $\Pi$ induces a partition of $M_f$ into at most $2^D$ monochromatic rectangles.
	\end{lemma}
	
	Given a protocol, its \emph{transcript} is the sequence of bits communicated.
	
	\begin{lemma}
	    For every transcript $z$ of some communication protocol, the set of inputs $(x,y)$ that generate $z$ is a rectangle.
	\end{lemma}
	
%
%
%
	
%
%
%

    Below, we recount the definitions of two multiparty communication models used in this work, namely the number-on-forehead and the number-in-hand models.
	
	\begin{definition}[``Number-on-forehead'' communication model; informal]
	    In the $k$-party ``number-on-forehead'' communication model, 
%
	    there are $k$ players and $k$ strings $x_1,\dots,x_k\in\bool^{n/k}$ and player $i$ gets all the strings except for $x_i$. The players are interested in computing a value $f\cprn{x_1,\dots,x_k}$, where $f\colon\bool^{n}\to\bool$ is some fixed function. 
%
	    We denote by $D^{(k)}(f)$ the number of bits that must be exchanged by the best possible number on forehead protocol solving $f$.
    \end{definition}
        
%
    
    We also use the following weaker communication model.
    
    \begin{definition}[``Number-in-hand'' communication model; informal]
	    In the $k$-party ``number-in-hand'' communication model,
%
	    there are $k$ players and $k$ strings $x_1,\dots,x_k\in\bool^{n/k}$ and player $i$ gets only $x_i$.
%
	    The players are interested in computing a value $f\cprn{x_1,\dots,x_k}$, where $f\colon\bool^{n}\to\bool$ is some fixed function.
%
	    We denote by $D^{(k\text{-}\mathsf{NIH})}(f)$ the number of bits that must be exchanged by the best possible communication protocol.
    \end{definition}
    
    Note that $D^{(k\text{-}\mathsf{NIH})}(f)\leq\prn{1-1/k}\cdot n+1$, for any $n$-variate Boolean function $f$, as if $k-1$ players write on the blackboard their string, then the player that did not reveal her input may compute $f\cprn{x_1,\dots,x_k}$ on her own and then publish it.
    
    For the communication models mentioned above, there are also bounded-error randomized versions, denoted by $R_\delta$, $R_\delta^{(k)}$, and $R^{(k\text{-}\mathsf{NIH})}_\delta$, respectively, where $0<\delta<1$ is an upper bound on the error probability of the protocol. In this setting, the players have access to some shared random string, say $r$, and the aforementioned error probability of the protocol is considered over the possible choices of $r$. Moreover, we require the error to be at most $\delta$ on each fixed choice of inputs.
    
    We can extend the definitions of the communication complexity measures, defined above, to classes of Boolean functions, in a natural way. That is, for any communication complexity measure $M\in\brkts{D,D^{(k)},D^{(k\text{-}\mathsf{NIH})},R_\delta,R_\delta^{(k)},R_\delta^{(k\text{-}\mathsf{NIH})}}$ and for any class of Boolean functions $\mathcal{G}$, we may define
    \[
        M\cprn{\mathcal{G}}:=\max_{g\in\mathcal{G}}M\cprn{g}.
    \]
	We note that throughout this paper, we denote by $n$ the number of input bits for the function regardless the communication models. In the $k$-party communication setting (either NOF or NIH), we assume without loss of generality that $n$ is divisible by $k$.
	
